\documentclass{article}
\usepackage{enumitem}
\usepackage{amsmath}
\begin{document}

\title{Chapter 32: Microbes in Human Welfare}
\date{}
\maketitle

\begin{enumerate}[label=Q\arabic*.]

\item The antibiotic penicillin was discovered by Alexander Fleming from:
\\(A) \textit{Streptomyces griseus}
\\(B) \textit{Penicillium notatum} (now \textit{P. chrysogenum})
\\(C) \textit{Bacillus subtilis}
\\(D) \textit{Aspergillus niger}

\item Streptomycin, the first antibiotic effective against tuberculosis, is produced by:
\\(A) \textit{Penicillium}
\\(B) \textit{Streptomyces griseus}
\\(C) \textit{Aspergillus}
\\(D) \textit{Bacillus}

\item Biogas (primarily methane) production in biogas plants involves:
\\(A) Aerobic decomposition of organic matter
\\(B) Anaerobic digestion by methanogenic archaea (e.g., \textit{Methanobacterium})
\\(C) Fermentation by yeast
\\(D) Photosynthetic bacteria

\item The BOD (Biological Oxygen Demand) of water is used to assess:
\\(A) The concentration of dissolved oxygen
\\(B) The amount of organic pollutants; higher BOD indicates more pollution
\\(C) The pH of water
\\(D) The presence of toxic chemicals

\item Assertion: Secondary sewage treatment uses microbes to reduce BOD of wastewater.\\
Reason: Aerobic microbes in activated sludge consume organic matter, reducing its concentration in wastewater.
\\(A) Both true, Reason is correct explanation
\\(B) Both true, Reason is not correct explanation
\\(C) Assertion true, Reason false
\\(D) Assertion false, Reason true

\item \textit{Trichoderma} is used as a biological control agent because it:
\\(A) Fixes nitrogen in soil
\\(B) Produces antibiotics that suppress fungal plant pathogens
\\(C) Decomposes cellulose for biogas production
\\(D) Solubilizes phosphate in soil

\item Mycorrhizal associations help plants by:
\\(A) Fixing atmospheric nitrogen
\\(B) Extending the root system for increased absorption of water and minerals (especially phosphorus)
\\(C) Producing phytohormones for growth
\\(D) Protecting roots from fungal diseases

\item Swiss cheese has characteristic holes (eyes) due to production of CO$_2$ by:
\\(A) \textit{Lactobacillus}
\\(B) \textit{Propionibacterium shermanii}
\\(C) \textit{Streptococcus thermophilus}
\\(D) \textit{Aspergillus}

\item Flocs in activated sludge treatment are:
\\(A) Inorganic precipitates settling in primary treatment
\\(B) Masses of bacteria held together by fungal filaments that digest organic matter in wastewater
\\(C) Chemical coagulants added to water
\\(D) Algal mats in the aeration tank

\item Biocontrol using \textit{Bacillus thuringiensis} (Bt) works because:
\\(A) Bt produces antibiotics that kill insects
\\(B) Bt produces Cry proteins (endotoxins) that form pores in the midgut of insect larvae when activated at alkaline pH
\\(C) Bt is a predatory bacterium that attacks insects
\\(D) Bt produces pheromones that disrupt insect mating

\item The primary treatment in sewage treatment plants involves:
\\(A) Biological digestion by microbes
\\(B) Physical removal of large solids and grit by screening and sedimentation
\\(C) Chemical disinfection with chlorine
\\(D) Anaerobic digestion in biogas digesters

\item Which microorganism is used in production of citric acid industrially?
\\(A) \textit{Saccharomyces cerevisiae}
\\(B) \textit{Aspergillus niger}
\\(C) \textit{Lactobacillus}
\\(D) \textit{Rhizobium}

\item \textit{Rhizobium} forms symbiotic association with legumes. The nitrogen fixed is directly available to the plant as:
\\(A) Nitrate (NO$_3^-$)
\\(B) Ammonia (NH$_3$) / ammonium (NH$_4^+$) incorporated into amino acids
\\(C) Atmospheric N$_2$ directly
\\(D) Nitrite (NO$_2^-$)

\item Match the following microbes with their applications:\\
1. \textit{Saccharomyces cerevisiae} $\rightarrow$ (a) Bread making and ethanol fermentation\\
2. \textit{Lactobacillus} $\rightarrow$ (b) Curd/yogurt production\\
3. \textit{Penicillium} $\rightarrow$ (c) Antibiotic production\\
4. \textit{Aspergillus niger} $\rightarrow$ (d) Citric acid production
\\(A) 1-a, 2-b, 3-c, 4-d
\\(B) 1-b, 2-a, 3-d, 4-c
\\(C) 1-c, 2-d, 3-a, 4-b
\\(D) 1-d, 2-c, 3-b, 4-a

\item The effluent from secondary sewage treatment is often treated with chlorine before release to:
\\(A) Remove remaining BOD
\\(B) Kill residual pathogenic microorganisms (disinfection)
\\(C) Remove heavy metals
\\(D) Neutralize the pH

\end{enumerate}
\end{document}
